% Daniel Gregorio-Torres - Currículum general en español
% Esta es la versión maestra en español. Copie este archivo antes de
% adaptarlo a una oferta de trabajo específica.

\documentclass[letterpaper,10pt]{article}

\usepackage[top=0.32in,bottom=0.32in,left=0.45in,right=0.45in]{geometry}
\usepackage[utf8]{inputenc}
\usepackage[T1]{fontenc}
\usepackage[spanish,es-nodecimaldot]{babel}
\usepackage{charter}
\usepackage{enumitem}
\usepackage[hidelinks]{hyperref}
\usepackage{tabularx}
\usepackage{xcolor}

\hypersetup{
  pdftitle={Currículum de Daniel Gregorio-Torres},
  pdfauthor={Daniel Gregorio-Torres},
  pdfsubject={Soporte de TI, administración de sistemas y desarrollo},
  pdfkeywords={soporte de TI, administración de sistemas, automatización, PowerShell, desarrollo web}
}

\pagestyle{empty}
\setlength{\parindent}{0pt}
\setlength{\tabcolsep}{0pt}
\linespread{0.98}
\setlist[itemize]{
  leftmargin=0.17in,
  label=\textbullet,
  itemsep=1pt,
  topsep=2pt,
  parsep=0pt,
  partopsep=0pt
}
\urlstyle{same}

\newcommand{\resumesection}[1]{%
  \vspace{4pt}
  {\large\bfseries\MakeUppercase{#1}}\par
  \vspace{2pt}
  \hrule
  \vspace{2.5pt}
}

\newcommand{\role}[4]{%
  \textbf{#1} \hfill \textbf{#2}\\
  \textit{#3} \hfill \textit{#4}\vspace{2.5pt}
}

\newcommand{\project}[3]{%
  \textbf{#1} \hfill \textit{#2}\\
  #3
}

\begin{document}

\begin{center}
  {\LARGE\bfseries Daniel Gregorio-Torres}\\[4pt]
  Beaverton, Oregón \textbar{}
  \href{mailto:djgregoriotorres@gmail.com}{djgregoriotorres@gmail.com}
  \\[2pt]
  \textbf{Sitio web:} \href{https://daniel.gregoriotorres.com/es/}{daniel.gregoriotorres.com} \textbar{}
  \textbf{GitHub:} \href{https://github.com/daniel1wnl}{daniel1wnl} \textbar{}
  \textbf{LinkedIn:} \href{https://linkedin.com/in/danielgregoriotorres}{danielgregoriotorres}
\end{center}

\resumesection{Resumen profesional}

Especialista en soporte de TI y desarrollador de sistemas con Licenciatura en Ciencias de la Computación y experiencia en soporte a usuarios, equipos, identidad, Microsoft 365 y entornos Windows híbridos. Desarrollo aplicaciones internas y automatizaciones con PowerShell que sustituyen procesos manuales y hacen que el trabajo de TI sea más repetible. Combino soporte técnico, desarrollo full-stack, documentación e implementaciones orientadas a la seguridad.

\resumesection{Habilidades técnicas}

\textbf{Operaciones de TI:} Windows, macOS, Linux, Microsoft 365, preparación de equipos, resolución de problemas, aplicación de parches, soporte de VPN, dispositivos móviles, documentación técnica\\
\textbf{Identidad y administración:} Active Directory, Microsoft Entra ID, Group Policy, Exchange, SSO/OAuth, permisos, ciclo de vida de usuarios\\
\textbf{Automatización y herramientas internas:} PowerShell, Power Apps, Power Automate, SharePoint Lists, JSON, Pester, Zendesk, IIS\\
\textbf{Desarrollo y datos:} Python, JavaScript, TypeScript, C\#, SQL, React, Node.js, ASP.NET Core, Flask, MongoDB, SQLite, APIs REST\\
\textbf{Infraestructura y entrega:} Git, GitHub, Docker, Google Cloud Run, Azure, AWS, monitoreo, respaldos, autohospedaje

\resumesection{Experiencia profesional}

\role{Especialista en soporte de TI}{Junio de 2024 -- Actualidad}{Streimer}{Portland, Oregón}
\begin{itemize}
  \item Brindo soporte de nivel 1 y nivel 2 para equipos Windows y dispositivos móviles, Microsoft 365, conectividad VPN, acceso de usuarios, hardware, software y tecnología operativa.
  \item Administro cuentas, permisos, equipos y acceso basado en políticas mediante Active Directory, Microsoft Entra ID, Exchange, Group Policy y herramientas de administración de dispositivos.
  \item Diseño automatizaciones con PowerShell para preparar estaciones de trabajo, gestionar el ciclo de vida de empleados, inventariar equipos, aplicar parches, realizar respaldos, validar procesos y estandarizar la resolución de problemas con registros y configuraciones seguras.
  \item Desarrollé y desplegué de forma independiente un sistema de inventario con Power Apps, SharePoint Lists y Power Automate que redujo el tiempo habitual para localizar materiales de aproximadamente 25--30 minutos a 5--10 minutos.
  \item Desarrollé SafetyCheck, una aplicación móvil con React que sustituyó las inspecciones en papel por registros validados, consultables y auditables, con búsqueda de equipos mediante códigos de barras.
  \item Creé un portal de intranet con ASP.NET Core alojado en IIS para consultar procedimientos mediante códigos QR, visualizar archivos PDF en dispositivos móviles y registrar confirmaciones estructuradas; documenté su despliegue, mantenimiento y soporte para tabletas.
\end{itemize}

\vspace{1pt}
\role{Proyecto patrocinado por un cliente}{Enero de 2023 -- Junio de 2023}{iGrafx / Proyecto final de Portland State University}{Remoto}
\begin{itemize}
  \item Colaboré en un equipo de seis personas para desarrollar un panel de lanzamientos con React, Node.js, Express y MongoDB que integraba datos de Jira, GitLab y SonarQube.
  \item Diseñé e implementé la página con la lista general de proyectos de la empresa y la página con los proyectos específicos de cada usuario.
  \item Contribuí con pruebas unitarias y documentación para mejorar la confiabilidad, el mantenimiento y la entrega del proyecto.
\end{itemize}

\resumesection{Proyectos técnicos seleccionados}

\project{Laboratorio doméstico}{Enero de 2026 -- Actualidad}{%
Mantengo un entorno privado para servicios en contenedores, almacenamiento, respaldos, monitoreo, acceso remoto seguro, automatización, resolución de problemas y documentación de recuperación. Practico la infraestructura como un ciclo de vida que abarca implementación, validación, mantenimiento y restauración.}

\vspace{2pt}
\project{Clasificaciones de Riot Games - Google Cloud Run}{Septiembre de 2022 -- Diciembre de 2022}{%
Desarrollé de forma independiente una aplicación con Flask que transforma datos en vivo de la API de Riot Games en clasificaciones de Valorant y League of Legends; la probé con Pytest, la empaqueté con Docker y la desplegué mediante Google Cloud Build y Cloud Run.}

\resumesection{Educación y desarrollo profesional}

\textbf{Licenciatura en Ciencias de la Computación, Minor en Administración de Empresas}
\hfill Junio de 2023\\
Portland State University \hfill Portland, Oregón\\[2pt]
\textbf{Portland State University - Web Development Program}
\hfill Septiembre de 2022\\
\textit{Credenciales verificadas:}
\href{https://www.credential.net/d5c287ef-5859-4de5-81e1-039377d2f609}{Finalización del programa} \textbar{}
\href{https://www.credential.net/97ce4aa2-99d3-4365-9dbf-cbbb423643b4}{JavaScript y APIs} \textbar{}
\href{https://www.credential.net/557f8a88-302f-4e70-bff1-672f41ed74d5}{HTML, CSS y Bootstrap}

\end{document}

